\appendix
\chapter{User Manual and Installation Instructions}

This appendix provides a step-by-step technical guide for configuring, building, and running SmartLib on a fresh development machine.

\section{Prerequisites Checklist}
\begin{itemize}
    \item \textbf{Java SE Development Kit (JDK 17 or higher):} Ensure \texttt{JAVA\_HOME} environment variable points to your JDK directory and \texttt{java -version} reports version 17 or above.
    \item \textbf{MySQL Community Server (Version 8.0 or higher):} Installed and running on default TCP port \texttt{3306}.
    \item \textbf{Git:} For cloning the repository.
    \item \textbf{Maven:} Optional. The repository includes an executable Maven Wrapper (\texttt{mvnw} for Linux/macOS, \texttt{mvnw.cmd} for Windows).
\end{itemize}

\section{Database Initialization}
\begin{enumerate}
    \item Open MySQL Command Line Client or MySQL Workbench:
    \begin{lstlisting}[language=SQL]
    CREATE DATABASE smartlib;
    USE smartlib;
    \end{lstlisting}
    \item Execute the database schema definition located at \texttt{src/main/resources/sql/schema.sql} to generate all 7 relational tables.
    \item Execute the sample data population script located at \texttt{src/main/resources/sample\_data.sql} to seed default categories, books, student records, and librarian administrative accounts.
\end{enumerate}

\section{Database Configuration}
\begin{enumerate}
    \item Navigate to \texttt{src/main/resources/}.
    \item Create a copy of \texttt{database.properties.example} and save it as \texttt{database.properties}.
    \item Update the database credentials to match your local MySQL configuration:
    \begin{lstlisting}
    db.url=jdbc:mysql://localhost:3306/smartlib
    db.user=root
    db.password=YOUR_MYSQL_PASSWORD
    \end{lstlisting}
\end{enumerate}

\section{Building and Executing the Application}
Open a terminal in the root repository folder:
\begin{enumerate}
    \item \textbf{Compile the Application:}
    \begin{lstlisting}[language=bash]
    # On Linux / macOS:
    ./mvnw compile
    
    # On Windows:
    .\mvnw.cmd compile
    \end{lstlisting}
    Verify that the build completes with zero errors.
    
    \item \textbf{Launch the GUI Desktop Application:}
    \begin{lstlisting}[language=bash]
    # On Linux / macOS:
    ./mvnw exec:java -Dexec.mainClass="com.college.library.gui.MainApplication"
    
    # On Windows:
    .\mvnw.cmd exec:java "-Dexec.mainClass=com.college.library.gui.MainApplication"
    \end{lstlisting}
\end{enumerate}

\clearpage
\section{Default User Credentials}
\begin{table}[h!]
    \centering
    \small
    \caption{Default Seed Credentials for Verification}
    \begin{tabularx}{\textwidth}{>{\raggedright\arraybackslash}p{3.2cm} >{\raggedright\arraybackslash}p{3.5cm} >{\raggedright\arraybackslash}p{3.2cm} >{\raggedright\arraybackslash}X}
        \toprule
        \textbf{Portal / Role} & \textbf{Identifier / Username} & \textbf{Default Password} & \textbf{Access Level} \\
        \midrule
        Librarian Portal & \texttt{admin} & \texttt{password123} & Full Administrative Access \\
        Student Portal & \texttt{TVE23CS001} & \texttt{password123} & Student Self-Service \\
        Student Portal & \texttt{TVE23CS002} & \texttt{password123} & Student Self-Service \\
        \bottomrule
    \end{tabularx}
    \label{tab:seed_creds}
\end{table}

\section{Troubleshooting Matrix}
\begin{table}[h!]
    \centering
    \small
    \caption{Common Setup Errors and Solutions}
    \begin{tabularx}{\textwidth}{>{\raggedright\arraybackslash}p{4.2cm} >{\raggedright\arraybackslash}p{4.0cm} >{\raggedright\arraybackslash}X}
        \toprule
        \textbf{Error Message} & \textbf{Root Cause} & \textbf{Corrective Action} \\
        \midrule
        \texttt{Communications link failure} & MySQL service is stopped or port is not 3306 & Verify that MySQL service is running (\texttt{systemctl status mysql} or \texttt{services.msc}). \\
        \midrule
        \texttt{Access denied for user 'root'} & Incorrect password in configuration & Verify and update root password in \texttt{src/main/}\allowbreak\texttt{resources/}\allowbreak\texttt{database.properties}. \\
        \midrule
        \texttt{Unknown database 'smartlib'} & Database not initialized & Run \texttt{CREATE DATABASE smartlib;} in MySQL Workbench. \\
        \midrule
        \texttt{Table 'smartlib.books' doesn't exist} & Schema script not executed & Import and execute \texttt{src/main/}\allowbreak\texttt{resources/sql/}\allowbreak\texttt{schema.sql}. \\
        \bottomrule
    \end{tabularx}
    \label{tab:troubleshooting}
\end{table}

\newpage
\chapter{Selected Source Code Listings}

\section{Generic Base Data Access Object: \texttt{CrudDAO.java}}
\begin{lstlisting}[language=Java, caption={Generic CrudDAO Interface}]
package com.college.library.dao;

import java.sql.SQLException;
import java.util.List;
import java.util.Optional;

public interface CrudDAO<T, ID> {
    void save(T entity) throws SQLException;
    Optional<T> findById(ID id) throws SQLException;
    List<T> findAll() throws SQLException;
    void update(T entity) throws SQLException;
    void delete(ID id) throws SQLException;
}
\end{lstlisting}

\section{Circulation Desk and Fine Engine: \texttt{BorrowService.java}}
\begin{lstlisting}[language=Java, caption={Core Methods in BorrowService.java}]
package com.college.library.service;

import com.college.library.dao.*;
import com.college.library.dto.ActiveBorrowingDTO;
import com.college.library.exception.LibraryException;
import com.college.library.model.*;

import java.math.BigDecimal;
import java.sql.SQLException;
import java.time.Clock;
import java.time.LocalDate;
import java.time.temporal.ChronoUnit;
import java.util.ArrayList;
import java.util.List;
import java.util.Optional;

public class BorrowService {
    private final BorrowTransactionDAO transactionDAO;
    private final BookDAO bookDAO;
    private final StudentDAO studentDAO;
    private final ReservationDAO reservationDAO;
    private final BorrowRequestDAO borrowRequestDAO;
    private final Clock clock;
    
    public static final BigDecimal DAILY_FINE_RATE = BigDecimal.valueOf(10.00);

    public BorrowService(BorrowTransactionDAO txDAO, BookDAO bookDAO, 
                         StudentDAO studentDAO, ReservationDAO resDAO, 
                         BorrowRequestDAO reqDAO) {
        this(txDAO, bookDAO, studentDAO, resDAO, reqDAO, Clock.systemDefaultZone());
    }

    public BorrowService(BorrowTransactionDAO txDAO, BookDAO bookDAO, 
                         StudentDAO studentDAO, ReservationDAO resDAO, 
                         BorrowRequestDAO reqDAO, Clock clock) {
        this.transactionDAO = txDAO;
        this.bookDAO = bookDAO;
        this.studentDAO = studentDAO;
        this.reservationDAO = resDAO;
        this.borrowRequestDAO = reqDAO;
        this.clock = clock;
    }

    public void issueBook(String studentId, String bookId, int librarianId) 
            throws LibraryException {
        try {
            Optional<Student> student = studentDAO.findById(studentId);
            if (student.isEmpty()) throw new LibraryException("Student not found.");
            if (!"ACTIVE".equalsIgnoreCase(student.get().getStatus())) {
                throw new LibraryException("Student account is not active.");
            }

            Optional<Book> book = bookDAO.findById(bookId);
            if (book.isEmpty()) throw new LibraryException("Book not found.");
            if (book.get().getStatus() != BookStatus.AVAILABLE) {
                throw new LibraryException("Book is currently unavailable (" + book.get().getStatus() + ").");
            }

            LocalDate issueDate = LocalDate.now(clock);
            LocalDate dueDate = issueDate.plusDays(14);

            BorrowTransaction tx = new BorrowTransaction();
            tx.setStudentId(studentId);
            tx.setBookId(bookId);
            tx.setLibrarianId(librarianId);
            tx.setIssueDate(issueDate);
            tx.setDueDate(dueDate);
            tx.setFine(BigDecimal.ZERO);
            tx.setStatus("ACTIVE");

            transactionDAO.save(tx);
            bookDAO.updateStatus(bookId, BookStatus.ISSUED);
        } catch (SQLException e) {
            throw new LibraryException("Database error during book issue: " + e.getMessage(), e);
        }
    }

    public boolean returnBook(String bookId, int librarianId) throws LibraryException {
        try {
            Optional<BorrowTransaction> activeTx = transactionDAO.findActiveByBookId(bookId);
            if (activeTx.isEmpty()) {
                throw new LibraryException("No active borrow transaction found for book ID: " + bookId);
            }

            BorrowTransaction tx = activeTx.get();
            LocalDate returnDate = LocalDate.now(clock);
            BigDecimal fine = calculateOverdueFine(tx.getDueDate());

            tx.setReturnDate(returnDate);
            tx.setFine(fine);
            tx.setStatus("COMPLETED");

            transactionDAO.update(tx);
            bookDAO.updateStatus(bookId, BookStatus.AVAILABLE);

            // Check reservation queue
            Optional<Book> bookOpt = bookDAO.findById(bookId);
            if (bookOpt.isPresent()) {
                Optional<Reservation> pendingRes = reservationDAO.findEarliestPendingByIsbn(bookOpt.get().getIsbn());
                if (pendingRes.isPresent()) {
                    Reservation r = pendingRes.get();
                    r.setStatus("FULFILLED");
                    reservationDAO.update(r);
                    return true; // Alert GUI that reservation is fulfilled
                }
            }
            return false;
        } catch (SQLException e) {
            throw new LibraryException("Database error during return: " + e.getMessage(), e);
        }
    }

    public BigDecimal calculateOverdueFine(LocalDate dueDate) {
        LocalDate today = LocalDate.now(clock);
        if (today.isAfter(dueDate)) {
            long daysOverdue = ChronoUnit.DAYS.between(dueDate, today);
            return BigDecimal.valueOf(daysOverdue).multiply(DAILY_FINE_RATE);
        }
        return BigDecimal.ZERO;
    }
}
\end{lstlisting}

\section{Centralized Design System: \texttt{UITheme.java}}
\begin{lstlisting}[language=Java, caption={Centralized UI Theme Styling Engine}]
package com.college.library.gui;

import javax.swing.*;
import javax.swing.border.EmptyBorder;
import javax.swing.table.DefaultTableCellRenderer;
import javax.swing.table.JTableHeader;
import java.awt.*;

public class UITheme {
    public static final Color BG_APP = new Color(248, 250, 252);
    public static final Color BG_CARD = Color.WHITE;
    public static final Color PRIMARY = new Color(37, 99, 235);
    public static final Color PRIMARY_HOVER = new Color(29, 78, 216);
    public static final Color TEXT_DARK = new Color(15, 23, 42);
    public static final Color TEXT_MUTED = new Color(100, 116, 139);
    public static final Color BORDER = new Color(226, 232, 240);

    public static final Font FONT_HEADER = new Font("Segoe UI", Font.BOLD, 22);
    public static final Font FONT_TITLE = new Font("Segoe UI", Font.BOLD, 15);
    public static final Font FONT_BODY = new Font("Segoe UI", Font.PLAIN, 13);
    public static final Font FONT_BOLD = new Font("Segoe UI", Font.BOLD, 13);

    public static JButton createPrimaryButton(String text) {
        JButton btn = new JButton(text);
        btn.setFont(FONT_BOLD);
        btn.setForeground(Color.WHITE);
        btn.setBackground(PRIMARY);
        btn.setFocusPainted(false);
        btn.setBorder(new EmptyBorder(8, 16, 8, 16));
        btn.setCursor(new Cursor(Cursor.HAND_CURSOR));
        return btn;
    }

    public static void styleTable(JTable table) {
        table.setFont(FONT_BODY);
        table.setRowHeight(36);
        table.setShowGrid(false);
        table.setIntercellSpacing(new Dimension(0, 0));
        table.setSelectionBackground(new Color(239, 246, 255));
        table.setSelectionForeground(TEXT_DARK);
        
        JTableHeader header = table.getTableHeader();
        header.setFont(FONT_BOLD);
        header.setBackground(new Color(241, 245, 249));
        header.setForeground(TEXT_DARK);
        header.setBorder(BorderFactory.createMatteBorder(0, 0, 1, 0, BORDER));
    }
}
\end{lstlisting}
